% add references as per sample relevant to your domain
\begin{thebibliography}{99}

\bibitem{ref1}
C. Phua, V. Lee, K. Smith-Miles, and R. Gayler, ``A comprehensive survey of data mining-based fraud detection research,'' \textit{arXiv preprint arXiv:1009.6119}, 2010.

\bibitem{ref2}
S. Bhattacharyya, S. Jha, K. Tharakunnel, and J. C. Westland, ``Data mining for credit card fraud: A comparative study,'' \textit{Decision Support Systems}, vol. 50, no. 3, pp. 602--613, 2011.

\bibitem{ref3}
A. Dal Pozzolo, O. Caelen, R. A. Johnson, and G. Bontempi, ``Calibrating probability with undersampling for unbalanced classification,'' in \textit{Proc. IEEE Symposium Series on Computational Intelligence}, 2015, pp. 159--166.

\bibitem{ref4}
A. C. Bahnsen, D. Aouada, A. Stojanovic, and B. Ottersten, ``Feature engineering strategies for credit card fraud detection,'' \textit{Expert Systems with Applications}, vol. 51, pp. 134--142, 2016.

\bibitem{ref5}
F. Carcillo, Y.-A. Le Borgne, O. Caelen, and G. Bontempi, ``Streaming active learning strategies for real-life credit card fraud detection: Assessment and visualization,'' \textit{International Journal of Data Science and Analytics}, vol. 5, no. 4, pp. 285--300, 2018.

\bibitem{ref6}
J. Jurgovsky, M. Granitzer, K. Ziegler, S. Calabretto, P.-E. Portier, L. Y. He-Guelton, and O. Caelen, ``Sequence classification for credit-card fraud detection,'' \textit{Expert Systems with Applications}, vol. 100, pp. 234--245, 2018.

\bibitem{ref7}
N. V. Chawla, K. W. Bowyer, L. O. Hall, and W. P. Kegelmeyer, ``SMOTE: Synthetic minority over-sampling technique,'' \textit{Journal of Artificial Intelligence Research}, vol. 16, pp. 321--357, 2002.

\bibitem{ref8}
H. He and E. A. Garcia, ``Learning from imbalanced data,'' \textit{IEEE Transactions on Knowledge and Data Engineering}, vol. 21, no. 9, pp. 1263--1284, 2009.

\bibitem{ref9}
T. Saito and M. Rehmsmeier, ``The precision-recall plot is more informative than the ROC plot when evaluating binary classifiers on imbalanced datasets,'' \textit{PLOS ONE}, vol. 10, no. 3, e0118432, 2015.

\bibitem{ref10}
M. Buda, A. Maki, and M. A. Mazurowski, ``A systematic study of the class imbalance problem in convolutional neural networks,'' \textit{Neural Networks}, vol. 106, pp. 249--259, 2018.

\bibitem{ref11}
T. Chen and C. Guestrin, ``XGBoost: A scalable tree boosting system,'' in \textit{Proc. 22nd ACM SIGKDD International Conference on Knowledge Discovery and Data Mining}, 2016, pp. 785--794.

\bibitem{ref12}
B. McMahan, E. Moore, D. Ramage, S. Hampson, and B. Ag\"uera y Arcas, ``Communication-efficient learning of deep networks from decentralized data,'' in \textit{Proc. 20th International Conference on Artificial Intelligence and Statistics}, PMLR, vol. 54, 2017, pp. 1273--1282.

\bibitem{ref13}
J. Kone\v{c}n\'y, H. B. McMahan, F. X. Yu, P. Richt\'arik, A. T. Suresh, and D. Bacon, ``Federated learning: Strategies for improving communication efficiency,'' \textit{arXiv preprint arXiv:1610.05492}, 2016.

\bibitem{ref14}
P. Kairouz et al., ``Advances and open problems in federated learning,'' \textit{Foundations and Trends in Machine Learning}, vol. 14, no. 1--2, pp. 1--210, 2021.

\bibitem{ref15}
Q. Yang, Y. Liu, T. Chen, and Y. Tong, ``Federated machine learning: Concept and applications,'' \textit{ACM Transactions on Intelligent Systems and Technology}, vol. 10, no. 2, pp. 1--19, 2019.

\bibitem{ref16}
K. Bonawitz, H. Eichner, W. Grieskamp, D. Huba, A. Ingerman, V. Ivanov, C. Kiddon, J. Kone\v{c}n\'y, S. Mazzocchi, H. B. McMahan, T. Van Overveldt, D. Petrou, D. Ramage, and J. Roselander, ``Towards federated learning at scale: System design,'' in \textit{Proc. MLSys}, 2019.

\bibitem{ref17}
T. Li, A. K. Sahu, M. Zaheer, M. Sanjabi, A. Talwalkar, and V. Smith, ``Federated optimization in heterogeneous networks,'' in \textit{Proc. MLSys}, 2020.

\bibitem{ref18}
K. Bonawitz, V. Ivanov, B. Kreuter, A. Marcedone, H. B. McMahan, S. Patel, D. Ramage, A. Segal, and K. Seth, ``Practical secure aggregation for privacy-preserving machine learning,'' in \textit{Proc. ACM SIGSAC Conference on Computer and Communications Security}, 2017, pp. 1175--1191.

\bibitem{ref19}
R. C. Geyer, T. Klein, and M. Nabi, ``Differentially private federated learning: A client level perspective,'' \textit{arXiv preprint arXiv:1712.07557}, 2017.

\bibitem{ref20}
L. Zhu, Z. Liu, and S. Han, ``Deep leakage from gradients,'' in \textit{Advances in Neural Information Processing Systems}, vol. 32, 2019.

\bibitem{ref21}
E. Bagdasaryan, A. Veit, Y. Hua, D. Estrin, and V. Shmatikov, ``How to backdoor federated learning,'' in \textit{Proc. 23rd International Conference on Artificial Intelligence and Statistics}, PMLR, vol. 108, 2020, pp. 2938--2948.

\bibitem{ref22}
Y. Liu, Y. Kang, C. Xing, T. Chen, and Q. Yang, ``A secure federated transfer learning framework,'' \textit{IEEE Intelligent Systems}, vol. 35, no. 4, pp. 70--82, 2020.

\bibitem{ref23}
N. K. Le, Y. Liu, Q. M. Nguyen, Q. Liu, F. Liu, Q. Cai, and S. Hirche, ``FedXGBoost: Privacy-preserving XGBoost for federated learning,'' \textit{arXiv preprint arXiv:2106.10662}, 2021.

\bibitem{ref24}
S. M. Lundberg and S.-I. Lee, ``A unified approach to interpreting model predictions,'' in \textit{Advances in Neural Information Processing Systems}, vol. 30, 2017.

\bibitem{ref25}
M. T. Ribeiro, S. Singh, and C. Guestrin, ``Why should I trust you? Explaining the predictions of any classifier,'' in \textit{Proc. 22nd ACM SIGKDD International Conference on Knowledge Discovery and Data Mining}, 2016, pp. 1135--1144.

\end{thebibliography}
